% !TEX root = ../main.tex
% Appendix F source fragment. Do not input this file into the six-page paper.
% The overall portfolio wrapper must load booktabs and hyperref, create the
% blank Appendix F title page, establish F.1/F.2/... page numbering, and define
% the optional content-page hooks below.
\providecommand{\PortfolioAppendixPageStyle}{}
\providecommand{\PortfolioContentPageCommand}{}

% The repository tables use the full page width for readable paths.
\onecolumn
\PortfolioAppendixPageStyle
\PortfolioContentPageCommand
\phantomsection
\addcontentsline{toc}{subsection}{F.1 Repository Overview}
\subsection*{F.1 Repository Overview}
\label{app:f-overview}

The project repository accompanies the electronic portfolio as a separate Loop
submission. The intended archive name is
\nolinkurl{A00049113_project_repository.zip}. It contains the implementation,
executed notebooks, tests and compact result records needed to inspect the
portfolio claims. The archive does not depend on an external cloud link.

The repository root is the directory that contains \nolinkurl{README.md},
\nolinkurl{sparam-surrogate/} and \nolinkurl{project_portfolio/}. Start with the
root README. It explains the directory structure, software environment, data
dependency and quickest inspection route. Table~\ref{tab:appendix-f-core}
identifies nine core items for detailed review.

\begin{table}[!htbp]
    \caption{Core repository items for detailed review.}
    \label{tab:appendix-f-core}
    \centering
    \scriptsize
    \begin{tabular}{@{}p{0.04\textwidth}p{0.34\textwidth}
                         p{0.54\textwidth}@{}}
        \toprule
        No. & Path or indexed file set & Purpose \\
        \midrule
        1
          & \nolinkurl{README.md}
          & Main repository guide. It gives setup, inspection and data-access
            instructions and repeats this core/supplementary distinction. \\
        2
          & \nolinkurl{project_portfolio/}
          & Six-page paper source, Appendix B--F source, references, figures,
            and the compact evidence used by the paper and appendices. Build
            products remain under its ignored \nolinkurl{build/} directory. \\
        3
          & \nolinkurl{sparam-surrogate/src/sparam_surrogate/}
          & Production Python package. It implements data alignment and loading,
            model wrappers, the reciprocal complete-matrix representation,
            physical diagnostics, run records and plotting utilities. \\
        4
          & \nolinkurl{sparam-surrogate/environment.yml},
            \nolinkurl{pyproject.toml}, \nolinkurl{configs/default.json},
            \nolinkurl{env_setup.sh}, and \nolinkurl{Makefile}
          & Environment, dependency, dataset, split, model and notebook-build
            configuration. These files define the seed and the expected
            directory layout. \\
        5
          & \nolinkurl{sparam-surrogate/notebooks/nb01_*.py} through
            \nolinkurl{nb08_*.py}, with their paired executed
            \nolinkurl{.ipynb} files
          & Complete exploratory and modelling path. The readable Jupytext
            \nolinkurl{.py} files are paired sources. The executed
            \nolinkurl{.ipynb} files retain displayed results. Individual
            filenames and purposes are listed in Table~\ref{tab:appendix-f-nbs}.
            \\
        6
          & \nolinkurl{sparam-surrogate/outputs/models/selected.json} and the
            compact records from the eight selected run directories
          & Frozen model selection and provenance. The compact record for each
            run comprises configuration, environment, manifest, metadata,
            metrics, training history and key figures. Large model binaries are
            treated separately in Section F.4. \\
        7
          & \nolinkurl{project_portfolio/evidence/}
          & Machine-readable result tables, selected-run provenance, bootstrap
            comparisons, physical diagnostics, metric-reproduction checks and
            the Appendix D neural-model inventory. \\
        8
          & \nolinkurl{sparam-surrogate/tests/}
          & Unit and integration tests for data handling, model contracts,
            physical reconstruction, run records and result utilities. \\
        9
          & \nolinkurl{sparam-surrogate/reports/pdf/nb01_*.pdf} through
            \nolinkurl{nb07_*.pdf}
          & Read-only notebook reports for assessors who prefer rendered PDFs.
            NB08 is available as an executed notebook and its exported diagrams
            are included under \nolinkurl{project_portfolio/media/appendix_d/}.
            \\
        \bottomrule
    \end{tabular}
\end{table}

\clearpage
\PortfolioContentPageCommand
\phantomsection
\addcontentsline{toc}{subsection}{F.2 Notebook Identifier Map}
\subsection*{F.2 Notebook Identifier Map}
\label{app:f-notebooks}

In the authored appendices, labels NB01 through NB08 are evidence locators for
the numbered Jupyter notebooks. They are not technical acronyms. Each notebook
has an executed \nolinkurl{.ipynb} file and a same-stem Jupytext
\nolinkurl{.py} source file. A range such as NB03--NB06 includes every notebook
from NB03 to NB06.
Paths in Table~\ref{tab:appendix-f-nbs} are relative to
\nolinkurl{sparam-surrogate/}.

\begin{table}[!htbp]
    \caption{Notebook labels, filenames and purposes.}
    \label{tab:appendix-f-nbs}
    \centering
    \scriptsize
    \begin{tabular}{@{}p{0.07\textwidth}p{0.40\textwidth}
                         p{0.44\textwidth}@{}}
        \toprule
        Label & Executed notebook filename & Purpose \\
        \midrule
        NB01
          & \nolinkurl{notebooks/nb01_dataset_exploration.ipynb}
          & Explores the chosen PCB dataset, parameter distributions, geometry
            relationships and representative Touchstone responses. \\
        NB02
          & \nolinkurl{notebooks/nb02_data_preprocessing.ipynb}
          & Aligns parameter and Touchstone records, creates the fixed
            design-level split, builds preprocessing artifacts and checks for
            split leakage. \\
        NB03
          & \nolinkurl{notebooks/nb03_non_neural_modelling.ipynb}
          & Trains and evaluates Scalar Ridge, Vector Ridge, Polynomial Ridge
            and Random Forest baselines. \\
        NB04
          & \nolinkurl{notebooks/nb04_neural_baseline.ipynb}
          & Trains and evaluates the point-wise neural and polynomial-neural
            multilayer-perceptron baselines. \\
        NB05
          & \nolinkurl{notebooks/nb05_curve_neural_model.ipynb}
          & Develops the whole-curve neural model and its validation-based
            frequency-feature, decoder and loss comparisons. \\
        NB06
          & \nolinkurl{notebooks/nb06_full_smatrix_physics.ipynb}
          & Trains the reciprocal complete complex S-matrix model and computes
            its accuracy and physical diagnostics. \\
        NB07
          & \nolinkurl{notebooks/nb07_selected_models_evaluation_analysis.ipynb}
          & Reloads the selected runs without retraining. It reproduces metrics,
            compares models and exports the final paper and Appendix E evidence.
            \\
        NB08
          & \nolinkurl{notebooks/nb08_appendix_d_model_graphs.ipynb}
          & Reloads selected estimators without retraining. It explains their
            structures and exports neural topology diagrams and the Appendix D
            model inventory. \\
        \bottomrule
    \end{tabular}
\end{table}

\clearpage
\PortfolioContentPageCommand
\phantomsection
\addcontentsline{toc}{subsection}{F.3 Inspection and Use Instructions}
\subsection*{F.3 Inspection and Use Instructions}
\label{app:f-instructions}

The core results can be inspected without retraining a model. A reviewer can
follow this short route:

\begin{enumerate}
    \item Open \nolinkurl{README.md}, then inspect the paper and Appendices D--E
          under \nolinkurl{project_portfolio/}.
    \item Open the executed NB07 notebook or its PDF report for the selected
          comparisons. Check the corresponding comma-separated value files in
          \nolinkurl{project_portfolio/evidence/} for machine-readable values.
    \item Open NB08 and the exported images in
          \nolinkurl{project_portfolio/media/appendix_d/} for model structures.
    \item Inspect \nolinkurl{outputs/models/selected.json} and the referenced
          compact run records to trace a selected result to its configuration,
          metrics and saved-run identifier.
\end{enumerate}

To run the code, create the declared Conda environment on a compatible Apple
Silicon system, then install the package from the repository root:

\begin{verbatim}
conda env create -f sparam-surrogate/environment.yml
conda run -n meng pip install -e ./sparam-surrogate
cd sparam-surrogate
conda run -n meng python -c \
  "import numpy.typing,pytest;raise SystemExit(pytest.main())"
\end{verbatim}

The neural dependencies in \nolinkurl{environment.yml} target Apple Silicon.
Other platforms require the appropriate TensorFlow build. The raw
SI/PI-Database dataset is not embedded in the core archive because of its size.
Its exact identifier and expected local path are recorded in
\nolinkurl{sparam-surrogate/configs/default.json}. After placing the dataset at
that path, the notebooks should be run in numerical order. NB07 and NB08 require
the selected-run records and any model binaries needed for regeneration.
The explicit \nolinkurl{numpy.typing} import is a compatibility step for the
installed NumPy/scikit-rf combination. With that step, all 251 tests pass in the
declared environment.

To verify a notebook pair without executing it, run:

\begin{verbatim}
conda run -n meng jupytext --diff \
  sparam-surrogate/notebooks/<notebook>.py \
  sparam-surrogate/notebooks/<notebook>.ipynb
\end{verbatim}

\clearpage
\PortfolioContentPageCommand
\phantomsection
\addcontentsline{toc}{subsection}{F.4 Supplementary Material and Packaging}
\subsection*{F.4 Supplementary Material and Packaging}
\label{app:f-supplementary}

The following material is supplementary unless it is explicitly included in a
core item above:

\begin{itemize}
    \item the raw dataset and generated processed-data caches;
    \item unselected exploratory runs, temporary predictions, routine logs and
          notebook execution stamps;
    \item large saved model binaries, including the selected Random Forest
          binary, which is approximately 6.1~GB;
    \item extra plots, archived reports and development notes that are not
          cited by the portfolio; and
    \item LaTeX intermediate files under \nolinkurl{project_portfolio/build/}.
\end{itemize}

The selected model binaries are useful for exact notebook regeneration, but
they are not needed to inspect the submitted source, executed notebooks,
compact run records or exported evidence. If Loop capacity permits, the final
archive may include the selected binaries as supplementary material. Otherwise,
their omission must remain explicit in the root README. In that case, the
workflow should be described as documented and version-controlled rather than
fully reproducible from the archive alone.

The Loop ZIP must be assembled from the working repository rather than from
\texttt{git archive}. Executed notebooks, rendered notebook reports, data and
run outputs are intentionally excluded from Git. Before submission, the archive
must be unpacked into a temporary directory. All nine core items must then be
opened or checked from that copy. External links may provide background or raw
data access, but they do not replace the source, result summaries and evidence
submitted inside the ZIP.
